% ============================================================================
%  Chapter 2 --- Context: Host Organisation, Industrial Context, ADN
% ============================================================================
\chapter{Project Context}
\label{ch:context}

\section{Host Organisation: Huawei Tunisia}

Huawei Technologies Co., Ltd. is a global leader in information and communications technology (ICT) infrastructure and smart devices. Founded in 1987 in Shenzhen, China, Huawei serves over three billion people worldwide and operates in more than 170 countries. In the telecommunications domain, Huawei provides end-to-end solutions spanning radio access networks (RAN), core networks, transport, and cloud infrastructure.

\textbf{Huawei Tunisia} operates as a regional entity within the North Africa \& Mediterranean business unit. The Cloud IT / Sales-Solution department --- the host of this internship --- focuses on positioning Huawei Cloud Stack (HCS) and associated AI/analytics solutions for Tunisian operators, principally Tunisie Telecom. The team bridges the gap between Huawei's global product portfolio (SmartCare, ADN, HCS) and the specific operational needs of the local market.

This project was conducted within this department, under the direct supervision of the Huawei Cloud IT team, with access to both Huawei's technical documentation and real production data from Tunisie Telecom.

\section{Tunisie Telecom: The Operator Context}

Tunisie Telecom (TT) is the incumbent telecommunications operator in Tunisia, providing fixed and mobile services to millions of subscribers. The Tunisian mobile market is characterised by:

\begin{itemize}
  \item A \textbf{prepaid-dominant} subscriber base ($\sim$80\% prepaid, $\sim$20\% postpaid), consistent with INTT 2023 market statistics;
  \item Three major operators: Tunisie Telecom, Ooredoo Tunisia, and Orange Tunisia;
  \item A growing 4G/LTE footprint with early 5G planning;
  \item Subscriber ARPU (Average Revenue Per User) in the range of 5--105 TND/month, reflecting a wide spectrum from low-recharge prepaid users to premium postpaid subscribers.
\end{itemize}

The data used in this project originates from TT's production environment, anonymised and provided through the Huawei partnership. This real-world grounding is a critical differentiator of the project.

\section{The Telecom OSS/BSS Landscape}

\subsection{Operations Support Systems (OSS)}

OSS encompasses all systems that manage the \emph{network} side of a telecommunications operator. Within Huawei's ecosystem, the OSS layer includes:

\begin{itemize}
  \item \textbf{Access:} RAN (Radio Access Network), FTTx (Fibre to the x), IP access
  \item \textbf{NOM:} Network Operations Management --- fault management, performance management, configuration
  \item \textbf{Core:} Core network elements including IoT platforms
  \item \textbf{Cloud:} Cloud infrastructure management
  \item \textbf{VAS:} Value Added Services management
\end{itemize}

OSS data manifests as KPIs measured per cell, per time window: throughput (Mbps), latency (ms), packet loss (\%), active user count, signal strength (RSRP in dBm), and hundreds of derived indicators.

\subsection{Business Support Systems (BSS)}

BSS encompasses all systems that manage the \emph{business} side: subscriber management, billing, revenue assurance, customer care. Within the context defined by the supervisor, the key BSS metrics include:

\begin{itemize}
  \item \textbf{Norm User:} Standard subscriber behavioural metrics (usage patterns, plan type, line type)
  \item \textbf{EAP:} Experience Analytics Platform --- experience-weighted analytics per subscriber
  \item \textbf{APPU:} Average Purchase Per User --- revenue per recharge or transaction event (TND)
  \item \textbf{DOU:} Data of Use --- monthly data consumption per subscriber (GB)
\end{itemize}

\subsection{The O+B Convergence Challenge: From Scattered Data to Unified Intelligence}

The fundamental challenge in modern telecom operations is \textbf{O+B Convergence}: the unification of OSS and BSS data to move beyond surface-level network analysis. As illustrated in Huawei's unified data architecture framework, the current state of many operators, including TT, is typically characterised by scattered data silos.

In this ``Now'' state, data is split across traditional, separate data warehouses for BSS, OSS, and B2B services. This fragmentation results in complex integration overhead and delayed analytics, ultimately creating an ``AI Blindness'' bottleneck. In this environment, AI initiatives remain stagnant, limited to task-specific chatbots or localized network optimizations that cannot see the broader business impact.

Our project directly implements Huawei's prescribed ``ToBe'' vision: transforming scattered data into centralized, AI-ready intelligence. By substituting siloed warehouses with a unified, three-layer data lake (raw, processed, curated), we establish the foundation for O+B convergence. This centralized data strategy eliminates the delay in analytics and powers an \emph{AI Development Platform} capable of answering complex, cross-domain questions autonomously:

\begin{quote}
\emph{``When cell X degrades, which subscribers are affected, what is the revenue impact, and what retention action should the CVM system trigger?''}
\end{quote}

Today, this question requires manual analysis across multiple systems. Our project automates this correlation, turning scattered AI experiments into scaled, unified productivity.

% ── Figure: Huawei NMS/OSS/BSS Stack ─────────────────────────────────────────
\begin{figure}[H]
\centering
\begin{tikzpicture}[
  node distance=0.4cm and 0.5cm,
  ossbox/.style={rectangle, rounded corners=3pt, minimum width=2cm, minimum height=0.8cm,
    draw=ossblue!70, fill=ossblue!12, font=\scriptsize\bfseries, align=center},
  bssbox/.style={rectangle, rounded corners=3pt, minimum width=2cm, minimum height=0.8cm,
    draw=bssgreen!70, fill=bssgreen!12, font=\scriptsize\bfseries, align=center},
  bigbox/.style={rectangle, rounded corners=5pt, minimum width=3.5cm, minimum height=1.2cm,
    draw=#1!60, fill=#1!10, font=\small\bfseries, align=center, line width=0.8pt},
]

  % OSS Layer
  \node[ossbox] (access) {Access\\(RAN, FTTx, IP)};
  \node[ossbox, right=of access] (nom) {NOM};
  \node[ossbox, right=of nom] (core) {Core\\(IoT)};
  \node[ossbox, right=of core] (cloud) {Cloud};
  \node[ossbox, right=of cloud] (vas) {VAS};
  
  % OSS label
  \node[above=0.3cm of nom, font=\small\bfseries, text=ossblue!80!black] (osslabel) {OSS Layer};
  
  \begin{scope}[on background layer]
    \node[fit=(access)(nom)(core)(cloud)(vas)(osslabel),
          draw=ossblue!40, fill=ossblue!4, rounded corners=6pt, inner sep=8pt] (ossgroup) {};
  \end{scope}

  % BSS Layer --- below
  \node[bssbox, below=1.5cm of access] (normuser) {Norm User};
  \node[bssbox, right=of normuser] (eap) {EAP};
  \node[bssbox, right=of eap] (appu) {APPU};
  \node[bssbox, right=of appu] (dou) {DOU};
  
  \node[above=0.3cm of eap, font=\small\bfseries, text=bssgreen!80!black] (bsslabel) {BSS Layer};
  
  \begin{scope}[on background layer]
    \node[fit=(normuser)(eap)(appu)(dou)(bsslabel),
          draw=bssgreen!40, fill=bssgreen!4, rounded corners=6pt, inner sep=8pt] (bssgroup) {};
  \end{scope}

  % NMS / CEM
  \node[bigbox=cemteal, above=1.2cm of ossgroup] (nms) {\textbf{NMS / CEM}\\SmartCare};
  
  % Data Lake
  \node[bigbox=datalake, below=1.2cm of bssgroup] (lake) {\textbf{Data Lake}\\(Huawei Analytics)};
  
  % AI Agent
  \node[bigbox=aiviolet, right=2cm of lake] (agent) {\textbf{AI Agent}\\(Our Project)};
  
  % CVM
  \node[bigbox=cvmorange, above=1.2cm of agent] (cvm) {\textbf{CVM}\\Business Output};

  % Arrows
  \draw[myarrow=cemteal!70!black] (nms) -- (ossgroup);
  \draw[myarrow=ossblue!70!black] (ossgroup) -- (lake);
  \draw[myarrow=bssgreen!70!black] (bssgroup) -- (lake);
  \draw[myarrow=datalake!70!black] (lake) -- (agent)
    node[arrlabel, midway, below] {O+B Data};
  \draw[myarrow=aiviolet!70!black] (agent) -- (cvm)
    node[arrlabel, midway, right] {Intelligence};
  
  % ADN label
  \node[below=0.3cm of lake, font=\footnotesize\itshape, text=huaweired!60!black] {
    ADN 5G/N3 --- Autonomous Driving Network
  };
  
\end{tikzpicture}
\caption{Huawei NMS/CEM/OSS/BSS ecosystem with data lake and AI Agent positioning (from supervisor notes).}
\label{fig:huawei-stack}
\end{figure}

\section{Customer Experience Management (CEM) and SmartCare}

Huawei SmartCare is a CEM platform deployed at operator sites worldwide. Its core functions include:

\begin{itemize}
  \item \textbf{KPI/KQI/CEI Computation:} From raw network counters, SmartCare derives Key Quality Indicators (KQI) and Customer Experience Indices (CEI) per subscriber, per service, per cell.
  \item \textbf{Demarcation:} When a service degrades, SmartCare identifies \emph{where} the problem is --- in the RAN, the core, the transport, or the service platform. This is the ``demarcation function.''
  \item \textbf{Experience Alerts:} Real-time alerts when a subscriber's experience drops below configurable thresholds.
\end{itemize}

Our platform maps directly to SmartCare's output:
\begin{itemize}
  \item Our \emph{SLA risk prediction} corresponds to SmartCare's experience risk scoring.
  \item Our \emph{anomaly detection} corresponds to SmartCare's demarcation function --- identifying abnormal cells and subscribers.
  \item Our \emph{OSS--BSS correlation engine} extends SmartCare's reach by linking experience degradation to revenue impact --- a function SmartCare itself does not perform.
\end{itemize}

\section{Customer Value Management (CVM) and the AI Intelligence Gap}

CVM systems consume intelligence about subscribers to drive business actions. However, without a unified data foundation bridging CEM and business data, the integration remains complex and reactive. A mature CVM demands:

\begin{itemize}
  \item \textbf{Churn prediction:} Identifying at-risk subscribers before they leave due to invisible network degradation.
  \item \textbf{Upsell triggers:} Recommending higher-tier plans to highly engaged users consuming large volumes of data (DoU).
  \item \textbf{Retention actions:} Proactive offers (bonus data, discounts) automatically targeted at subscribers experiencing documented service drops (SLA risks).
  \item \textbf{Revenue impact scoring:} Quantifying the exact business cost (via APPU and Revenue) of network anomalies.
\end{itemize}

As positioned in Huawei's target architecture, the consumers of a unified AI platform include AI assistants, marketing copilots, and intelligent network analysts. Our AI Agent fills this specific gap---the intelligence layer bridging CEM outputs to CVM inputs. It produces exactly the required intelligent triggers: churn-correlated anomaly flags, SLA risk scores, and revenue impact assessments, natively exposing these cross-domain insights via REST API for immediate business action.

\section{The Autonomous Driving Network (ADN) Paradigm: Achieving L4 Autonomy}

Huawei's ADN vision charts the evolution of telecom networks through progressive levels of autonomy, mirroring the automotive industry. Recently updated industry frameworks recognize the progression of AI intelligence from basic Chatbots (L1) and Reasoners (L2) to fully autonomous, complex organizations (L5). 

\begin{itemize}
  \item \textbf{L1 --- Assisted (Chatbot):} Able to answer basic questions; manual operations with simple tool support.
  \item \textbf{L2 --- Partial (Reasoner):} Able to think; some automated, human-supervised workflows.
  \item \textbf{L3 --- Conditional (Agent):} \emph{Able to do}. The era of AI Agents taking automated actions under human supervision.
  \item \textbf{L4 --- Highly Autonomous (Innovator):} \emph{Able to invent and decide}. AI manages complex workflows and executes automated remediation without human intervention.
  \item \textbf{L5 --- Fully Autonomous (Organization):} Able to manage large and complex systems; the network manages itself end-to-end.
\end{itemize}

Historically, network analytics have been stuck in passive analytics (L1/L2) or conditional alerts (L3). Our project is explicitly designed to pierce the \textbf{L4 (Highly Autonomous)} threshold. It achieves this by moving beyond merely generating dashboards. The pipeline acts as an autonomous intelligence engine that actively ingests O+B data, orchestrates parallel ML models (Gradient Boosting, Isolation Forests), and critically, \textbf{generates and pushes autonomous action triggers} (e.g., CVM retention payloads, automated SLA breach escalations) without human intervention. By proving that the agent can "decide and act" to prevent revenue loss, this project serves as a foundational L4 proof-of-concept for the ultimate L5 Huawei vision.

% ── Figure: ADN Levels ────────────────────────────────────────────────────────
\begin{figure}[H]
\centering
\begin{tikzpicture}[
  level/.style={
    rectangle, rounded corners=3pt, minimum width=10cm, minimum height=0.9cm,
    font=\small, align=center, line width=0.6pt,
  },
]
  \node[level, draw=gray!50, fill=gray!8] (l1) at (0,0)
    {\textbf{L1 --- Chatbot (Assisted):} Able to answer questions};
  \node[level, draw=gray!50, fill=gray!8] (l2) at (0,1.1)
    {\textbf{L2 --- Reasoner (Partial):} Able to think / assist workflows};
  \node[level, draw=gray!50, fill=gray!8] (l3) at (0,2.2)
    {\textbf{L3 --- Agent (Conditional):} Able to do / recommend};
  \node[level, draw=aiviolet!70, fill=aiviolet!15, line width=1.2pt] (l4) at (0,3.3)
    {\textbf{L4 --- Innovator (Highly Autonomous):} Able to invent / decide $\leftarrow$ \textit{Our Project}};
  \node[level, draw=gray!50, fill=gray!8] (l5) at (0,4.4)
    {\textbf{L5 --- Organization (Fully Autonomous):} Manages complex systems};
  
  % Arrow on the side
  \draw[-{Stealth[length=8pt]}, line width=1.5pt, aiviolet!60] (-5.8,0) -- (-5.8,4.4)
    node[midway, rotate=90, above, font=\small\bfseries, text=aiviolet!70!black] {Increasing Autonomy};
\end{tikzpicture}
\caption{Huawei ADN autonomy levels. Our project targets L4 --- Highly Autonomous.}
\label{fig:adn-levels}
\end{figure}

\section{Project Scope and Boundaries}

\subsection{In Scope}

\begin{itemize}
  \item Ingestion of real anonymised OSS/BSS data from Tunisie Telecom via Huawei
  \item Three-layer data lake: Raw $\to$ Processed $\to$ Curated (MinIO / OBS)
  \item Three ML models: SLA risk (GBR), OSS anomalies (IF), BSS anomalies (IF)
  \item OSS--BSS statistical correlation engine (Pearson/Spearman)
  \item RESTful API for insight consumption
  \item Docker Compose orchestration (5 services)
  \item Architectural mapping to HCS (OBS, RDS, ECS/CCE)
  \item Model evaluation with real data (precision, recall, F1, confusion matrix)
\end{itemize}

\subsection{Explicitly Out of Scope}

\begin{itemize}
  \item Vendor-grade telecom assurance at production scale (millions of subscribers)
  \item Real-time streaming ingestion (our batch pipeline demonstrates the architecture)
  \item Direct SmartCare API integration (we ingest CEM-equivalent data from exports)
  \item Production HCS deployment (architectural readiness, not production cutover)
  \item CVM action execution (the AI Agent produces intelligence; CVM acts on it)
\end{itemize}

\section{Chapter Summary}

This chapter established the industrial context of the project: the Huawei CEM/CVM ecosystem, the SmartCare platform, the O+B convergence challenge, and the ADN paradigm. The project is positioned as the intelligence layer between CEM and CVM, implemented cloud-natively and trained on real Tunisie Telecom data. The next chapter surveys the state of the art in AI-driven telecom operations, positioning our approach relative to existing academic and industrial work.
